\chapter{State of the Art}
\label{chap:state-of-the-art}

\begin{introduction}
This chapter reviews the background and existing work relevant to the viewer developed in this thesis. It states how the literature was gathered, then covers the application of digital twins, introduced in Chapter~\ref{chap:introduction}, to physical infrastructure, the data sources and models that infrastructure digital twins draw on, the shift from dedicated Geographic Information System (GIS) tooling to real-time game engines for geospatial visualisation, and the handling of simulated entities along with physical ones. It then examines existing digital twin platforms and viewers, and how tightly each couples its rendering to a particular backend and data model. The chapter closes by stating the gap this thesis targets.
\end{introduction}

\section{Review Methodology}
\label{sec:review-methodology}

The literature underlying this chapter, and the digital twin background in Chapter~\ref{chap:introduction}, was gathered through a keyword search across Scopus, IEEE Xplore, Google Scholar, and ScienceDirect, using broad term combinations such as ``digital twin'', ``digital twin'' AND ``game engine'', and ``digital twins in game engines''. Because ``digital twin'' spans disciplines from manufacturing to healthcare to civil engineering, these broad terms returned a large, heterogeneous candidate pool; the Scopus search alone returned 320 records. This pool was narrowed by title and abstract relevance to the intersection this thesis actually addresses, meaning infrastructure digital twins, real-time game-engine visualisation, and generalisable viewer architectures, discarding the substantial majority of results that fell in unrelated domains (manufacturing digital twins, medical imaging, structural digital twins with no visualisation component, and similar).

The directly relevant set that remained after this narrowing is small, reflecting how narrow this specific intersection is in the published literature rather than an artefact of the search. Some of the works ultimately cited, including the foundational digital twin papers by \textcite{grieves2016digitaltwin} and \textcite{glaessgen2012digitaltwin}, were identified directly as known origin points for the concept rather than surfaced by the keyword search; further sources, including \textcite{broo2021digitaltwins}, \textcite{kherbache2022ditto}, and the geospatial-visualisation and related-viewer-system literature in Sections~\ref{sec:geospatial-vis} and~\ref{sec:related-viewer-systems}, were found by snowballing: following the reference lists and citing works of the papers already found this way.

\section{Digital Twins for Physical Infrastructure}
\label{sec:dt-infrastructure}

\textcite{broo2021digitaltwins} narrow the digital twin concept, introduced in Section~\ref{sec:context}, to infrastructure specifically, describing digital twins in this domain as an emerging idea without a commonly accepted definition, but converging on a realistic digital representation of assets, processes, or systems in the built environment, kept synchronised with their physical counterpart. In the mobility sub-domain, entities are the components of a road network and the sensing and telematics equipment that instruments them, as described in Section~\ref{sec:context}, and digital twins are expected to support applications such as traffic management, predictive maintenance, incident detection, and scenario evaluation~\cite{fuller2020digitaltwin}.

Architecturally, an infrastructure digital twin is conventionally split into a backend that models and serves entity data and a viewer that renders it. The backend's role is to abstract over the heterogeneity of the physical layer, with its many device types and protocols, and expose a uniform interface. Eclipse Ditto is open-source middleware built for exactly this: it represents physical devices as digital twins behind a common HTTP and WebSocket API, decoupling client applications from how each device actually communicates~\cite{kherbache2022ditto}. Ditto is explicitly a data layer; it provides no visualisation. \textcite{kherbache2022ditto} pair it with Eclipse Hono for device connectivity and demonstrate the combination as a digital twin network for the industrial IoT. Still, any three-dimensional or geospatial representation of the twins it manages must be built separately.

This backend/viewer separation is near-universal in the surveyed literature: the system that models the twin's data and the system that renders it are treated, and built, as distinct concerns. Where systems fall short of that separation \emph{within} the viewer, meaning the viewer's own code is written against one backend and one entity schema, is the subject of Section~\ref{sec:related-viewer-systems}.

This coupling leaves a gap that is not simply anecdotal. \textcite{masoumi2023citydigitaltwins} systematically reviewed 41 published city digital twin studies and found that 90\% of them sit at an initial-to-medium maturity level, concentrated on 3D modelling, monitoring, and visualisation as opposed to genuine two-way interaction; they identify near-real-time data analytics and improved BIM/GIS integration as open needs rather than solved problems. This independently corroborates, across a much larger sample than this chapter surveys directly, that the field as a whole has not moved much past the coupled default described above.

\section{Data Sources and Models for Infrastructure Digital Twins}
\label{sec:dt-data}

An infrastructure digital twin draws on several kinds of data, each with its own structure and rate of change. Producing and storing that data is the backend's concern rather than the viewer's (Section~\ref{sec:scope}), but the kinds of data involved determine what the viewer has to display: semantically rich models of static assets, continuous streams of sensor readings, and large numbers of moving vehicles.

\subsection{Building Information Modelling}

Building Information Modelling (BIM) is a methodology for creating and managing digital representations of the physical and functional characteristics of built assets, with detailed geometric, semantic, and relational information about each element, such as its dimensions, materials, and lifecycle attributes~\cite{azhar2011bim}. Although BIM originated in the architecture, engineering, and construction sector, it is increasingly applied to large infrastructure such as roads, bridges, and tunnels. Within digital twin frameworks, BIM models provide a high-fidelity structural and semantic layer; enriched with operational and real-time data, a BIM model stops being a static design artefact and becomes an active part of a digital twin, supporting asset management, condition monitoring, and maintenance planning~\cite{deng2021bimdigitaltwins}. For a viewer, such models are the source of asset footprints and attributes, such as those in a road operator's inventory of signage, barriers, and pavement sections.

\subsection{Sensors and Data Acquisition}

Sensors measure properties of the environment or of a structure, such as traffic flow, vehicle speed, weather conditions, structural stress, or air quality. In intelligent transportation systems they are the primary interface between the physical infrastructure and digital platforms, providing continuous data from the road network~\cite{felemban2014its}. Common types include inductive loop detectors, cameras, radar, LiDAR, environmental sensors, and structural health monitoring devices. The data they produce is heterogeneous, frequent, and geographically distributed, which makes it hard to integrate, process, and visualise. For a viewer, each sensor is an entity whose state changes live, such as an inclinometer on a retaining wall or a roadside traffic counter.

\subsection{Connected Vehicles}

Connected vehicles exchange data with other vehicles (vehicle-to-vehicle, V2V), with roadside infrastructure (vehicle-to-infrastructure, V2I), and with networked systems (vehicle-to-network, V2N), collectively referred to as vehicle-to-everything (V2X) communication~\cite{abboud2016v2x}. They report their position, speed, and trajectory, as well as observed road conditions and traffic events, and integrating this information into a digital twin of a road corridor gives a detailed, up-to-date picture of traffic~\cite{wu2025v2xdigitaltwin}. For a viewer, connected vehicles mean large numbers of moving entities whose positions change with every update, similar to the simulated vehicles discussed in Section~\ref{sec:simulated-entities}.

\section{Geospatial Visualisation: From GIS to Game Engines}
\label{sec:geospatial-vis}

Rendering a digital twin of physical infrastructure at real-world scale requires placing content accurately on a geographic reference frame. This has historically been the province of GIS software~\cite{goodchild2010giscience} and of standardised 3D city models: CityGML and its JSON encoding, CityJSON, describe urban objects, such as buildings, roads, vegetation, and city furniture, with defined geometry and semantics, and most urban digital twin platforms are built around such a model. The bridge from a BIM-side model to this representation is itself an active area: \textcite{xu2019ifc3dtiles} convert Industry Foundation Classes (IFC) building models into the 3D Tiles format this thesis also streams through, and \textcite{lam2024digitaltwinsmartcity} similarly map IFC into CityGML and an OWL/RDF semantic graph, visualising the result through both Cesium Ion and Unreal Engine; both treat 3D Tiles and CityGML as a common target format for otherwise incompatible authoring tools rather than something specific to one project.

Real-time 3D game engines have become a practical alternative as they have gained geospatial capabilities. \textcite{rantanen2023opengeospatial} integrate open geospatial data into a game engine for urban digital twin applications and describe the data-conversion pipeline this requires; \textcite{ene2026urbandigitaltwins} survey the combined use of GIS and game engine systems for urban digital twins and find game engines increasingly viable where interactivity and rendering quality matter more than analytical GIS functionality. Both note the trade-off: game engines bring strong rendering and interaction but must be bridged to geospatial data and coordinate systems they do not natively understand. \textcite{virtanen2020pointclouds} push this further for raw sensor data, extending a game engine to interactively render dense coloured point clouds instead of only pre-processed city models.

The coordinate systems involved recur throughout this thesis and are worth stating. Geographic positions are almost always expressed in the World Geodetic System 1984 (WGS84), the global reference system used by GPS, which models the Earth as an ellipsoid and identifies each point by its latitude, longitude, and height above that ellipsoid~\cite{nima2000wgs84}. Latitude and longitude are angles rather than distances, so they cannot be used directly as positions in a game engine, which works in a flat Cartesian space measured in fixed units. A common intermediate step is the Earth-Centred, Earth-Fixed (ECEF) frame, a Cartesian system in metres whose origin is the Earth's centre of mass and which rotates with the Earth. A WGS84 position converts to ECEF with a fixed formula, and ECEF positions can then be expressed relative to a local origin near the area being viewed, which yields the flat, metre-scale coordinates an engine expects.

The bridge used in this work is Cesium for Unreal~\cite{cesiumforunreal}, a free and open-source plugin that embeds a full-scale WGS84 globe in Unreal Engine and streams terrain, imagery, and 3D content into it. The streaming format is 3D Tiles, an OGC Community Standard for organising and delivering massive heterogeneous 3D geospatial datasets, such as point clouds, photogrammetry, BIM and CAD models, and 3D buildings, as a spatially indexed hierarchy suitable for progressive loading~\cite{ogc3dtiles}. Cesium for Unreal was chosen over building geospatial placement from scratch, and over a pure web or GIS viewer, because it retains Unreal Engine's rendering and interaction capabilities, relevant for a viewer intended to generalise across infrastructure types with potentially rich per-entity visuals, while still anchoring content to accurate real-world coordinates. Section~\ref{sec:geospatial-placement} in Chapter~\ref{chap:architecture} describes how this thesis builds a single coordinate-conversion seam on top of Cesium's georeferencing rather than scattering coordinate mathematics across the application.

\section{Simulated Entities Along with Physical Ones}
\label{sec:simulated-entities}

A viewer for a mobility digital twin is increasingly expected to display not only physical infrastructure but also entities produced by simulation and streamed through the same backend. Two cases arose in this work. The first is traffic: microscopic traffic simulators such as SUMO model individual vehicles' movement through a road network~\cite{lopez2018sumo}, producing a continuous, high-churn stream of short-lived vehicle entities, each with a position that changes with every update. The second is hazard modelling: a wildfire-spread simulation produces a projected burn perimeter and spread cone that evolve and are attached to an ignition-point entity.

These entities differ from physical infrastructure in ways that matter to the viewer. They are numerous and ephemeral, update far more frequently, may share a single fleet-wide lifetime rather than expiring individually, and may have their geometry delivered as GeoJSON rather than as a mesh. A viewer whose entity handling is specialised to static infrastructure will not accommodate them without change. This reinforces the need, argued throughout this thesis, for an entity model that is generic over what an entity is rather than assuming one shape.

A related case from outside the mobility domain is \textcite{yang2023cesiumlandslide}'s web-based emergency-response system, also built on Cesium Digital Earth, which maps live IoT and sensor-network data for rescue workers, vehicles, and monitored landslide sites onto the globe for search-and-rescue coordination. It is a hazard-adjacent parallel to the wildfire case above, streaming heterogeneous physical entities (personnel, equipment, instrumented sites) through one viewer. However, the paper does not report whether its entity handling generalises beyond this fixed set, as this thesis argues it should.

\section{Digital Twin Platforms and Related Viewer Systems}
\label{sec:related-viewer-systems}

Digital twin platforms for cities and mobility are now numerous, and a recent systematic review maps how IoT, artificial intelligence, and digital twins are combined across smart-city systems~\cite{sacoto2025smartcities}. The systems discussed below are those closest to the viewer in this thesis: each pairs a source of data with a 3D or geospatial view of infrastructure or mobility, and together they are useful comparators for how tightly the view is coupled to the data behind it.

City2Twin~\cite{rafamatanantsoa2024city2twin} presents an open urban digital twin pipeline that integrates municipal open data into a 3D city model for visualisation and analysis. Like most urban digital twin platforms surveyed by \textcite{ene2026urbandigitaltwins}, it is organised around a specific data pipeline and city-model format, typically CityGML or CityJSON. It is not designed to be repointed at an arbitrary IoT backend or entity schema. The 3D city model is both its input and its internal representation.

Cesium for Unreal itself, and the wider Cesium platform, are backend-agnostic in principle: they render whatever geospatial data they are given. But they ship no entity model, no live-update mechanism, and no interaction layer for IoT-style digital twins. Applications built on Cesium, including those surveyed in Section~\ref{sec:geospatial-vis}, still implement their own data layer, and that layer is typically written for one project's backend and one project's entity types. The geospatial foundation is reusable; the digital twin logic on top of it generally is not.

The closest domain match in the surveyed literature is \textcite{lv2024highwaysandbox}'s digital sandbox platform, also built on Unreal Engine, for highway route planning and design. It integrates BIM data from a fixed commercial toolchain (Bentley OpenRoads Designer and OpenBridge Modeler) with a photogrammetry-derived GIS terrain model into what the authors call a digital base model, then builds route-selection, traffic-simulation, and component-query modules on top of it. Two differences from the viewer in this thesis are instructive: it is a design-time planning tool built from a static, pre-processed model rather than a viewer kept current by a live backend connection, and its component types, meaning bridges, tunnels, and interchanges, are fixed to that one BIM/GIS pipeline rather than expressed through a generic, extensible entity model. Even a viewer this close in domain and engine choice does not attempt either of the two axes of generalisation this thesis argues for.

CARLA~\cite{dosovitskiy2017carla}, an open-source urban driving simulator built on Unreal Engine, shows how far a game engine can be taken in the mobility domain: it provides high-fidelity 3D urban environments, realistic traffic scenarios, and simulated sensors such as cameras and LiDAR. Its purpose, however, is developing and testing autonomous-driving systems at the level of individual vehicles, not monitoring a real road network through a live connection to infrastructure data. Closer to the connected road corridors that mobility digital twins target, \textcite{wu2025v2xdigitaltwin} propose a digital twin framework for integrating the physical and virtual sides of V2X-enabled connected-vehicle corridors; its contribution is that integration framework, and it does not address a viewer that generalises across entity types or backends.

Outside the geospatial domain, \textcite{robles2023opentwins}'s OpenTwins is an open-source framework for \emph{compositional} digital twins, in which individual entities and subsystems are linked and orchestrated with 3D visualisation, IoT data streams, and machine-learning predictions. Its composition-over-monolith stance is close in spirit to the entity model in Chapter~\ref{chap:architecture}, and it demonstrates that open, generalisable digital twin tooling is achievable; its target domain, however, is Industry 4.0 process orchestration, not geospatial mobility infrastructure, and it does not address placing entities on a real-world globe or streaming them by geographic area.

The academic and open-source systems above have commercial, closed-source counterparts, and these are worth naming because they show the coupled default is not a limitation of small or under-resourced projects: it persists in mature, widely deployed enterprise tooling. Bentley's iTwin platform~\cite{bentleyitwin} is a cloud digital twin ecosystem built around \emph{iModels}, Bentley's own proprietary schema (the Base Infrastructure Schema, BIS), populated primarily from Bentley's own authoring tools, such as MicroStation, OpenRoads, and OpenBuildings, the same product family used for the BIM side of the highway sandbox platform above. Siemens' Xcelerator portfolio~\cite{siemensxcelerator} plays a similar role for industrial and infrastructure IoT, connecting a digital twin to sensor data through Siemens' own device and data-model ecosystem. Esri's ArcGIS platform~\cite{esriarcgis}, extended for city- and infrastructure-scale digital twins through products such as ArcGIS Urban, takes the same shape from the GIS side: entities and their attributes are modelled as Esri geodatabase feature classes, and the surrounding tooling, analysis, and visualisation are all built to that one format. In each case, the viewer, the data model, and often the authoring tools are a single vendor-integrated product; none is designed to point at an arbitrary third-party backend or accept an entity type it was not built to know about, which is precisely the coupling this thesis's architecture avoids.

Across these systems a common pattern appears: the viewer's rendering and interaction code is written against one backend's data model, one fixed set of entity types, or both. This is a defensible choice for a single-deployment viewer, but it means the viewer cannot be reused for a different infrastructure domain or backend without editing its core, because no seam separates ``what an entity looks like and does'' or ``where its data comes from'' from the rest of the application.

\section{Summary and Research Gaps}
\label{sec:research-gaps}

The literature establishes the digital twin as a physical entity, a virtual representation, and a live connection between them. It establishes infrastructure digital twins as a backend/viewer split in which the backend (for example, Eclipse Ditto) models entity data and the viewer renders it. It shows that game engines paired with geospatial plugins such as Cesium for Unreal, streaming standardised formats such as 3D Tiles, are now a viable alternative to GIS-specific tooling for building that viewer. It also shows that such viewers are increasingly expected to handle very different kinds of data at once: semantically rich asset models, live sensor readings, moving vehicles, and simulated entities such as traffic and hazards.

What the surveyed viewer systems do not generally address is generalisation and decoupling \emph{within} the viewer. Existing systems tend to fix both the entity types they can display and the backend they fetch from, so neither can change without changing the viewer's core code, and the viewer cannot run without a live instance of its specific backend. The gap matters most where one viewer must serve several infrastructure domains at once, since a viewer fixed to one of them cannot serve the others. This thesis targets that gap: it presents a viewer whose entity model and data-access layer are each a single replaceable seam rather than assumptions baked throughout the rendering and interaction code, as detailed in Chapter~\ref{chap:architecture}.
